\documentclass[11pt]{scrartcl}
\usepackage[a4paper,margin=2.6cm]{geometry}
\usepackage{../common/rlstudy}

\title{Extended reading --- Week 9: Proximal Policy Optimization}
\date{\today}
\begin{document}
\maketitle

\section*{1.\ Schulman, Wolski, Dhariwal, Radford, Klimov (2017) --- Proximal Policy Optimization Algorithms}
Preprint, \href{https://arxiv.org/abs/1707.06347}{arXiv:1707.06347}.

\paragraph{Summary.} PPO alternates interaction with surrogate optimization and enables repeated minibatch updates on a collected batch. The clipped and adaptive KL-penalty variants have favorable empirical comparisons on continuous-control and Atari tasks; these experiments are not convergence proofs.
\paragraph{Why it matters for this week.} Derive the four cases formed by advantage sign and ratio location relative to the clipping interval. Separate improvement incentives from penalties for harmful changes.
\paragraph{Annotated key passage (abstract).}
\begin{quote}\emph{``a novel objective function that enables multiple epochs of minibatch updates''}\end{quote}
Batch reuse explains why old log probabilities stay fixed. The clipped surrogate is pointwise no larger than the unclipped surrogate, because it takes their minimum. This is not a lower bound on actual policy return.


\section*{2.\ Huang, Dossa, Raffin, Kanervisto, Wang (2022) --- The 37 Implementation Details of Proximal Policy Optimization}
ICLR Blog Track, 25 March 2022. \href{https://iclr-blog-track.github.io/2022/03/25/ppo-implementation-details/}{Article and reference implementation}.

\paragraph{Summary.} The authors trace implementation revisions and reproduce results across classic control, Atari, robotics, recurrent policies, and multidiscrete actions. Their 37-item checklist contains task-specific choices; every item is not required for every environment.
\paragraph{Why it matters for this week.} Audit stored Gaussian samples, action clipping, normalization, and rollout boundaries. The course's detail numbers are local labels, not an exact copy of the article's ordering.
\paragraph{Annotated key passage (introduction).}
\begin{quote}\emph{``conclusions are in MuJoCo tasks and do not necessarily transfer to other games such as Atari''}\end{quote}
Evaluate choices in their intended task. An ablation label is meaningful only if disabling the option changes the implementation.


\section*{3.\ Engstrom, Ilyas, Santurkar, Tsipras, Janoos, Rudolph, Madry (2020) --- Implementation Matters in Deep Policy Gradients}
ICLR 2020, \href{https://arxiv.org/abs/2005.12729}{arXiv:2005.12729}.

\paragraph{Summary.} Controlled PPO/TRPO comparisons show that code-level optimizations affect both measured performance and the trust-region behavior attributed to the algorithm. Comparisons can change when implementation conditions change.
\paragraph{Why it matters for this week.} Match architectures, optimizers, interaction budgets, and evaluation protocols when comparing losses.
\paragraph{Annotated passage (paraphrase of the abstract).}
\begin{quote}Implementation details affect the performance and behavior attributed to a policy-gradient algorithm.\end{quote}
This is a paraphrase, not a direct quotation. The result motivates controlled comparisons; it does not establish that PPO and TRPO are interchangeable on every task or that one checklist eliminates their differences.


\section*{4.\ Andrychowicz, Raichuk, Stanczyk, Orsini, Girgin, Marinier, Hussenot, Geist, Pietquin, Michalski, Gelly, Bachem (2021) --- What Matters In On-Policy Reinforcement Learning? A Large-Scale Empirical Study}
ICLR 2021, \href{https://arxiv.org/abs/2006.05990}{arXiv:2006.05990}.

\paragraph{Summary.} This broad empirical study investigates design decisions in on-policy deep RL, covering objectives, architectures, optimization, and data collection. The effects are conditional on the tested environments and protocols.
\paragraph{Why it matters for this week.} Select an ablation from the study, then state its task, budget, and seeds before interpreting the effect.
\paragraph{Annotated passage (paraphrase of the abstract).}
\begin{quote}Empirical investigation of many design choices can inform on-policy RL implementations.\end{quote}
This is a paraphrase, not a direct quotation. Recommendations motivate measurements; they are not convergence theorems. There is no universal ordering of which implementation choice matters most.


\section*{5.\ Ouyang, Wu, Jiang, et al.\ (2022) --- Training language models to follow instructions with human feedback (InstructGPT)}
NeurIPS 2022, \href{https://arxiv.org/abs/2203.02155}{arXiv:2203.02155}.

\paragraph{Summary.} InstructGPT combines supervised demonstrations, a preference reward model, and PPO fine-tuning to improve instruction following under human evaluation. It is an application of PPO, not evidence that all preference-training systems use the same optimizer.
\paragraph{Why it matters for this week.} Distinguish the old rollout policy from the SFT reference used in KL reward regularization.
\paragraph{Annotated passage (paraphrase of the method).}
\begin{quote}Divergence from the supervised reference is penalized during reward optimization.\end{quote}
This is a paraphrase, not a direct quotation. A finite penalty discourages drift but does not eliminate reward hacking or guarantee safety. Week 18 supplies a small synthetic analogue, rather than a frontier-model reproduction.

\section*{Sourced from}
\texttt{research/deeprl\_core\_sota.md} \S 2.6 and \texttt{research/llm\_rl\_sota.md} (RLHF section).

\end{document}
